% =============================================================================
% EIF-Based Power Analysis for Mean and Difference in Means
% =============================================================================
% This section can be % =============================================================================
% EIF-Based Power Analysis for Mean and Difference in Means
% =============================================================================
% This section can be % =============================================================================
% EIF-Based Power Analysis for Mean and Difference in Means
% =============================================================================
% This section can be % =============================================================================
% EIF-Based Power Analysis for Mean and Difference in Means
% =============================================================================
% This section can be \input{eif_power} in main.tex when ready

\section{EIF-Based Power Analysis}
\label{sec:eif-power}

We derive power formulas from the perspective of semiparametric efficiency theory.
The efficient influence function (EIF) characterizes the asymptotically optimal
estimator for a given statistical model, and its variance directly determines
the power of hypothesis tests.

\subsection{EIF for the Mean with Auxiliary Predictions}

Consider estimating the population mean $\theta^\star = \E[Y]$ using:
\begin{itemize}
    \item Labeled data: $(X_i, Y_i, f_i)$ for $i = 1, \ldots, n$, where $f_i = f(X_i)$
    \item Unlabeled data: $(\widetilde{X}_j, \widetilde{f}_j)$ for $j = 1, \ldots, N$
\end{itemize}

\paragraph{The classical influence function.}
Without auxiliary predictions, the influence function for the mean is simply
\[
\psi_0(Y) = Y - \theta^\star.
\]
The sample mean $\bar{Y}_n$ is the corresponding M-estimator with asymptotic variance $\Var(Y)/n$.

\paragraph{The augmented influence function.}
With auxiliary predictions $f(X)$, we can construct an augmented influence function:
\begin{equation}\label{eq:eif-augmented}
\psi_\lambda(Y, f) = Y - \theta^\star + \lambda\big(f - \E[f]\big),
\end{equation}
where $\lambda \in \mathbb{R}$ is a tuning parameter.

The key observation is that $\E[\psi_\lambda(Y, f)] = 0$ for any $\lambda$ since
$\E[Y] = \theta^\star$ and $\E[f - \E[f]] = 0$. This means the estimating equation
\[
\frac{1}{n}\sum_{i=1}^n \psi_\lambda(Y_i, f_i) = 0
\]
yields an unbiased estimator of $\theta^\star$ for any choice of $\lambda$.

\paragraph{Variance of the augmented estimator.}
The variance of $\psi_\lambda(Y, f)$ is
\begin{align*}
\Var\big(\psi_\lambda(Y, f)\big)
&= \Var(Y) + \lambda^2 \Var(f) + 2\lambda\Cov(Y - \theta^\star, f - \E[f])\\
&= \sigma_Y^2 + \lambda^2 \sigma_f^2 + 2\lambda\Cov(Y, f).
\end{align*}

Minimizing over $\lambda$ gives the \emph{efficient} choice:
\begin{equation}\label{eq:eif-lambda-star}
\lambda^\star_{\text{EIF}} = -\frac{\Cov(Y, f)}{\sigma_f^2},
\end{equation}
with optimal variance
\begin{equation}\label{eq:eif-opt-var}
\Var\big(\psi_{\lambda^\star}(Y, f)\big)
= \sigma_Y^2 - \frac{\Cov(Y, f)^2}{\sigma_f^2}
= \sigma_Y^2 \big(1 - \rho^2_{Yf}\big),
\end{equation}
where $\rho_{Yf} = \Cov(Y, f)/(\sigma_Y \sigma_f)$ is the correlation between $Y$ and $f$.

\paragraph{Connection to PPI++.}
The PPI++ estimator~\eqref{eq:ppi++-mean-est} with optimal $\lambda^\star$ from~\eqref{eq:lambda-star-mean}
achieves exactly this efficient variance in~\eqref{eq:eif-opt-var}, confirming that
PPI++ is the semiparametrically efficient estimator for the mean when auxiliary predictions are available.

\subsection{Power Inversion for EIF-Based Inference}

Given the efficient variance~\eqref{eq:eif-opt-var}, we can derive power formulas by
the standard approach of inverting the power equation.

\paragraph{Setup.}
Consider testing
\[
H_0: \theta^\star = \theta_0 \quad \text{vs} \quad H_1: \theta^\star = \theta_0 + \Delta
\]
at significance level $\alpha$ with target power $1 - \beta$.

\paragraph{Asymptotic distribution under $H_1$.}
The efficient estimator $\widehat{\theta}_{\text{EIF}}$ satisfies
\[
\sqrt{n}\big(\widehat{\theta}_{\text{EIF}} - \theta^\star\big) \xrightarrow{d} \mathcal{N}\big(0, \sigma_{\text{EIF}}^2\big),
\]
where $\sigma_{\text{EIF}}^2 = \sigma_Y^2(1 - \rho_{Yf}^2)$ is the efficient variance.

\paragraph{Power formula.}
The power of the two-sided Wald test is
\begin{equation}\label{eq:eif-power}
\text{Power} = \Phi\left(-z_{1-\alpha/2} + \frac{|\Delta|}{\sigma_{\text{EIF}}/\sqrt{n}}\right)
+ \Phi\left(-z_{1-\alpha/2} - \frac{|\Delta|}{\sigma_{\text{EIF}}/\sqrt{n}}\right).
\end{equation}

\paragraph{Required sample size.}
Inverting~\eqref{eq:eif-power} for target power $1 - \beta$ gives
\begin{equation}\label{eq:eif-n-required}
n \ge \frac{\sigma_Y^2(1 - \rho_{Yf}^2)}{\big(\Delta / (z_{1-\alpha/2} + z_{1-\beta})\big)^2}.
\end{equation}

\paragraph{Efficiency gain.}
Compared to the classical sample mean (without predictions), which requires
\[
n_{\text{classical}} \ge \frac{\sigma_Y^2}{\big(\Delta / (z_{1-\alpha/2} + z_{1-\beta})\big)^2},
\]
the EIF-based approach reduces the required sample size by a factor of $(1 - \rho_{Yf}^2)$.
When $\rho_{Yf} = 0.9$, this is a $81\%$ reduction.

\subsection{EIF for Difference in Means}

For the two-sample setting with independent groups $A$ and $B$, the parameter of interest is
$\Delta^\star = \theta_A^\star - \theta_B^\star$.

\paragraph{Augmented influence function.}
The efficient influence function for the difference is
\[
\psi_{\lambda_A, \lambda_B}(Y_A, f_A, Y_B, f_B)
= \big(Y_A + \lambda_A(f_A - \E[f_A])\big) - \big(Y_B + \lambda_B(f_B - \E[f_B])\big).
\]

With group-specific optimal tuning:
\begin{equation}\label{eq:eif-ttest-lambda}
\lambda_A^\star = -\frac{\Cov(Y_A, f_A)}{\sigma_{f,A}^2}, \qquad
\lambda_B^\star = -\frac{\Cov(Y_B, f_B)}{\sigma_{f,B}^2},
\end{equation}
the efficient variance of the difference estimator is
\begin{equation}\label{eq:eif-ttest-var}
\Var(\widehat{\Delta}_{\text{EIF}})
= \frac{\sigma_{Y,A}^2(1 - \rho_A^2)}{n_A} + \frac{\sigma_{Y,B}^2(1 - \rho_B^2)}{n_B},
\end{equation}
where $\rho_A, \rho_B$ are the group-specific correlations between $Y$ and $f$.

\paragraph{Power formula for two-sample EIF test.}
Under balanced design ($n_A = n_B = n$) and equal variances ($\sigma_{Y,A} = \sigma_{Y,B} = \sigma_Y$),
the power is
\begin{equation}\label{eq:eif-ttest-power}
\text{Power} = \Phi\left(-z_{1-\alpha/2} + \frac{|\Delta|}{\sqrt{2}\sigma_Y\sqrt{(1-\bar\rho^2)/n}}\right)
+ \Phi\left(-z_{1-\alpha/2} - \frac{|\Delta|}{\sqrt{2}\sigma_Y\sqrt{(1-\bar\rho^2)/n}}\right),
\end{equation}
where $\bar\rho^2 = (\rho_A^2 + \rho_B^2)/2$ is the average squared correlation.

\paragraph{Required sample size per group.}
\begin{equation}\label{eq:eif-ttest-n}
n \ge \frac{2\sigma_Y^2(1 - \bar\rho^2)}{\big(\Delta / (z_{1-\alpha/2} + z_{1-\beta})\big)^2}.
\end{equation}

\subsection{Practical Calibration}

The EIF-based formulas require knowledge of:
\begin{itemize}
    \item $\sigma_Y^2$: estimated from labeled data as $\widehat{\sigma}_Y^2 = \frac{1}{n-1}\sum_{i=1}^n(Y_i - \bar{Y})^2$
    \item $\rho_{Yf}$: estimated from labeled data as the sample correlation between $Y_i$ and $f_i$
\end{itemize}

\paragraph{From MSE to correlation.}
For a well-calibrated predictor with $\E[f(X)] \approx \E[Y]$, the MSE relates to the correlation via
\[
\text{MSE} = \E[(Y - f)^2] = \sigma_Y^2 + \sigma_f^2 - 2\Cov(Y, f) = \sigma_Y^2(1 - \rho^2) + (\sigma_f - \rho\sigma_Y)^2.
\]
When $f$ is the best linear predictor of $Y$, we have $\sigma_f = \rho\sigma_Y$ and $\text{MSE} = \sigma_Y^2(1 - \rho^2)$.

\paragraph{From $R^2$ to sample size reduction.}
In regression settings, the $R^2$ of predicting $Y$ from $f$ equals $\rho_{Yf}^2$.
Thus, the sample size reduction factor is exactly $(1 - R^2)$.

\begin{tcolorbox}
\textbf{Rule of thumb:} An auxiliary predictor with $R^2 = 0.5$ (explaining 50\% of variance)
reduces the required labeled sample size by 50\%.
\end{tcolorbox}
 in main.tex when ready

\section{EIF-Based Power Analysis}
\label{sec:eif-power}

We derive power formulas from the perspective of semiparametric efficiency theory.
The efficient influence function (EIF) characterizes the asymptotically optimal
estimator for a given statistical model, and its variance directly determines
the power of hypothesis tests.

\subsection{EIF for the Mean with Auxiliary Predictions}

Consider estimating the population mean $\theta^\star = \E[Y]$ using:
\begin{itemize}
    \item Labeled data: $(X_i, Y_i, f_i)$ for $i = 1, \ldots, n$, where $f_i = f(X_i)$
    \item Unlabeled data: $(\widetilde{X}_j, \widetilde{f}_j)$ for $j = 1, \ldots, N$
\end{itemize}

\paragraph{The classical influence function.}
Without auxiliary predictions, the influence function for the mean is simply
\[
\psi_0(Y) = Y - \theta^\star.
\]
The sample mean $\bar{Y}_n$ is the corresponding M-estimator with asymptotic variance $\Var(Y)/n$.

\paragraph{The augmented influence function.}
With auxiliary predictions $f(X)$, we can construct an augmented influence function:
\begin{equation}\label{eq:eif-augmented}
\psi_\lambda(Y, f) = Y - \theta^\star + \lambda\big(f - \E[f]\big),
\end{equation}
where $\lambda \in \mathbb{R}$ is a tuning parameter.

The key observation is that $\E[\psi_\lambda(Y, f)] = 0$ for any $\lambda$ since
$\E[Y] = \theta^\star$ and $\E[f - \E[f]] = 0$. This means the estimating equation
\[
\frac{1}{n}\sum_{i=1}^n \psi_\lambda(Y_i, f_i) = 0
\]
yields an unbiased estimator of $\theta^\star$ for any choice of $\lambda$.

\paragraph{Variance of the augmented estimator.}
The variance of $\psi_\lambda(Y, f)$ is
\begin{align*}
\Var\big(\psi_\lambda(Y, f)\big)
&= \Var(Y) + \lambda^2 \Var(f) + 2\lambda\Cov(Y - \theta^\star, f - \E[f])\\
&= \sigma_Y^2 + \lambda^2 \sigma_f^2 + 2\lambda\Cov(Y, f).
\end{align*}

Minimizing over $\lambda$ gives the \emph{efficient} choice:
\begin{equation}\label{eq:eif-lambda-star}
\lambda^\star_{\text{EIF}} = -\frac{\Cov(Y, f)}{\sigma_f^2},
\end{equation}
with optimal variance
\begin{equation}\label{eq:eif-opt-var}
\Var\big(\psi_{\lambda^\star}(Y, f)\big)
= \sigma_Y^2 - \frac{\Cov(Y, f)^2}{\sigma_f^2}
= \sigma_Y^2 \big(1 - \rho^2_{Yf}\big),
\end{equation}
where $\rho_{Yf} = \Cov(Y, f)/(\sigma_Y \sigma_f)$ is the correlation between $Y$ and $f$.

\paragraph{Connection to PPI++.}
The PPI++ estimator~\eqref{eq:ppi++-mean-est} with optimal $\lambda^\star$ from~\eqref{eq:lambda-star-mean}
achieves exactly this efficient variance in~\eqref{eq:eif-opt-var}, confirming that
PPI++ is the semiparametrically efficient estimator for the mean when auxiliary predictions are available.

\subsection{Power Inversion for EIF-Based Inference}

Given the efficient variance~\eqref{eq:eif-opt-var}, we can derive power formulas by
the standard approach of inverting the power equation.

\paragraph{Setup.}
Consider testing
\[
H_0: \theta^\star = \theta_0 \quad \text{vs} \quad H_1: \theta^\star = \theta_0 + \Delta
\]
at significance level $\alpha$ with target power $1 - \beta$.

\paragraph{Asymptotic distribution under $H_1$.}
The efficient estimator $\widehat{\theta}_{\text{EIF}}$ satisfies
\[
\sqrt{n}\big(\widehat{\theta}_{\text{EIF}} - \theta^\star\big) \xrightarrow{d} \mathcal{N}\big(0, \sigma_{\text{EIF}}^2\big),
\]
where $\sigma_{\text{EIF}}^2 = \sigma_Y^2(1 - \rho_{Yf}^2)$ is the efficient variance.

\paragraph{Power formula.}
The power of the two-sided Wald test is
\begin{equation}\label{eq:eif-power}
\text{Power} = \Phi\left(-z_{1-\alpha/2} + \frac{|\Delta|}{\sigma_{\text{EIF}}/\sqrt{n}}\right)
+ \Phi\left(-z_{1-\alpha/2} - \frac{|\Delta|}{\sigma_{\text{EIF}}/\sqrt{n}}\right).
\end{equation}

\paragraph{Required sample size.}
Inverting~\eqref{eq:eif-power} for target power $1 - \beta$ gives
\begin{equation}\label{eq:eif-n-required}
n \ge \frac{\sigma_Y^2(1 - \rho_{Yf}^2)}{\big(\Delta / (z_{1-\alpha/2} + z_{1-\beta})\big)^2}.
\end{equation}

\paragraph{Efficiency gain.}
Compared to the classical sample mean (without predictions), which requires
\[
n_{\text{classical}} \ge \frac{\sigma_Y^2}{\big(\Delta / (z_{1-\alpha/2} + z_{1-\beta})\big)^2},
\]
the EIF-based approach reduces the required sample size by a factor of $(1 - \rho_{Yf}^2)$.
When $\rho_{Yf} = 0.9$, this is a $81\%$ reduction.

\subsection{EIF for Difference in Means}

For the two-sample setting with independent groups $A$ and $B$, the parameter of interest is
$\Delta^\star = \theta_A^\star - \theta_B^\star$.

\paragraph{Augmented influence function.}
The efficient influence function for the difference is
\[
\psi_{\lambda_A, \lambda_B}(Y_A, f_A, Y_B, f_B)
= \big(Y_A + \lambda_A(f_A - \E[f_A])\big) - \big(Y_B + \lambda_B(f_B - \E[f_B])\big).
\]

With group-specific optimal tuning:
\begin{equation}\label{eq:eif-ttest-lambda}
\lambda_A^\star = -\frac{\Cov(Y_A, f_A)}{\sigma_{f,A}^2}, \qquad
\lambda_B^\star = -\frac{\Cov(Y_B, f_B)}{\sigma_{f,B}^2},
\end{equation}
the efficient variance of the difference estimator is
\begin{equation}\label{eq:eif-ttest-var}
\Var(\widehat{\Delta}_{\text{EIF}})
= \frac{\sigma_{Y,A}^2(1 - \rho_A^2)}{n_A} + \frac{\sigma_{Y,B}^2(1 - \rho_B^2)}{n_B},
\end{equation}
where $\rho_A, \rho_B$ are the group-specific correlations between $Y$ and $f$.

\paragraph{Power formula for two-sample EIF test.}
Under balanced design ($n_A = n_B = n$) and equal variances ($\sigma_{Y,A} = \sigma_{Y,B} = \sigma_Y$),
the power is
\begin{equation}\label{eq:eif-ttest-power}
\text{Power} = \Phi\left(-z_{1-\alpha/2} + \frac{|\Delta|}{\sqrt{2}\sigma_Y\sqrt{(1-\bar\rho^2)/n}}\right)
+ \Phi\left(-z_{1-\alpha/2} - \frac{|\Delta|}{\sqrt{2}\sigma_Y\sqrt{(1-\bar\rho^2)/n}}\right),
\end{equation}
where $\bar\rho^2 = (\rho_A^2 + \rho_B^2)/2$ is the average squared correlation.

\paragraph{Required sample size per group.}
\begin{equation}\label{eq:eif-ttest-n}
n \ge \frac{2\sigma_Y^2(1 - \bar\rho^2)}{\big(\Delta / (z_{1-\alpha/2} + z_{1-\beta})\big)^2}.
\end{equation}

\subsection{Practical Calibration}

The EIF-based formulas require knowledge of:
\begin{itemize}
    \item $\sigma_Y^2$: estimated from labeled data as $\widehat{\sigma}_Y^2 = \frac{1}{n-1}\sum_{i=1}^n(Y_i - \bar{Y})^2$
    \item $\rho_{Yf}$: estimated from labeled data as the sample correlation between $Y_i$ and $f_i$
\end{itemize}

\paragraph{From MSE to correlation.}
For a well-calibrated predictor with $\E[f(X)] \approx \E[Y]$, the MSE relates to the correlation via
\[
\text{MSE} = \E[(Y - f)^2] = \sigma_Y^2 + \sigma_f^2 - 2\Cov(Y, f) = \sigma_Y^2(1 - \rho^2) + (\sigma_f - \rho\sigma_Y)^2.
\]
When $f$ is the best linear predictor of $Y$, we have $\sigma_f = \rho\sigma_Y$ and $\text{MSE} = \sigma_Y^2(1 - \rho^2)$.

\paragraph{From $R^2$ to sample size reduction.}
In regression settings, the $R^2$ of predicting $Y$ from $f$ equals $\rho_{Yf}^2$.
Thus, the sample size reduction factor is exactly $(1 - R^2)$.

\begin{tcolorbox}
\textbf{Rule of thumb:} An auxiliary predictor with $R^2 = 0.5$ (explaining 50\% of variance)
reduces the required labeled sample size by 50\%.
\end{tcolorbox}
 in main.tex when ready

\section{EIF-Based Power Analysis}
\label{sec:eif-power}

We derive power formulas from the perspective of semiparametric efficiency theory.
The efficient influence function (EIF) characterizes the asymptotically optimal
estimator for a given statistical model, and its variance directly determines
the power of hypothesis tests.

\subsection{EIF for the Mean with Auxiliary Predictions}

Consider estimating the population mean $\theta^\star = \E[Y]$ using:
\begin{itemize}
    \item Labeled data: $(X_i, Y_i, f_i)$ for $i = 1, \ldots, n$, where $f_i = f(X_i)$
    \item Unlabeled data: $(\widetilde{X}_j, \widetilde{f}_j)$ for $j = 1, \ldots, N$
\end{itemize}

\paragraph{The classical influence function.}
Without auxiliary predictions, the influence function for the mean is simply
\[
\psi_0(Y) = Y - \theta^\star.
\]
The sample mean $\bar{Y}_n$ is the corresponding M-estimator with asymptotic variance $\Var(Y)/n$.

\paragraph{The augmented influence function.}
With auxiliary predictions $f(X)$, we can construct an augmented influence function:
\begin{equation}\label{eq:eif-augmented}
\psi_\lambda(Y, f) = Y - \theta^\star + \lambda\big(f - \E[f]\big),
\end{equation}
where $\lambda \in \mathbb{R}$ is a tuning parameter.

The key observation is that $\E[\psi_\lambda(Y, f)] = 0$ for any $\lambda$ since
$\E[Y] = \theta^\star$ and $\E[f - \E[f]] = 0$. This means the estimating equation
\[
\frac{1}{n}\sum_{i=1}^n \psi_\lambda(Y_i, f_i) = 0
\]
yields an unbiased estimator of $\theta^\star$ for any choice of $\lambda$.

\paragraph{Variance of the augmented estimator.}
The variance of $\psi_\lambda(Y, f)$ is
\begin{align*}
\Var\big(\psi_\lambda(Y, f)\big)
&= \Var(Y) + \lambda^2 \Var(f) + 2\lambda\Cov(Y - \theta^\star, f - \E[f])\\
&= \sigma_Y^2 + \lambda^2 \sigma_f^2 + 2\lambda\Cov(Y, f).
\end{align*}

Minimizing over $\lambda$ gives the \emph{efficient} choice:
\begin{equation}\label{eq:eif-lambda-star}
\lambda^\star_{\text{EIF}} = -\frac{\Cov(Y, f)}{\sigma_f^2},
\end{equation}
with optimal variance
\begin{equation}\label{eq:eif-opt-var}
\Var\big(\psi_{\lambda^\star}(Y, f)\big)
= \sigma_Y^2 - \frac{\Cov(Y, f)^2}{\sigma_f^2}
= \sigma_Y^2 \big(1 - \rho^2_{Yf}\big),
\end{equation}
where $\rho_{Yf} = \Cov(Y, f)/(\sigma_Y \sigma_f)$ is the correlation between $Y$ and $f$.

\paragraph{Connection to PPI++.}
The PPI++ estimator~\eqref{eq:ppi++-mean-est} with optimal $\lambda^\star$ from~\eqref{eq:lambda-star-mean}
achieves exactly this efficient variance in~\eqref{eq:eif-opt-var}, confirming that
PPI++ is the semiparametrically efficient estimator for the mean when auxiliary predictions are available.

\subsection{Power Inversion for EIF-Based Inference}

Given the efficient variance~\eqref{eq:eif-opt-var}, we can derive power formulas by
the standard approach of inverting the power equation.

\paragraph{Setup.}
Consider testing
\[
H_0: \theta^\star = \theta_0 \quad \text{vs} \quad H_1: \theta^\star = \theta_0 + \Delta
\]
at significance level $\alpha$ with target power $1 - \beta$.

\paragraph{Asymptotic distribution under $H_1$.}
The efficient estimator $\widehat{\theta}_{\text{EIF}}$ satisfies
\[
\sqrt{n}\big(\widehat{\theta}_{\text{EIF}} - \theta^\star\big) \xrightarrow{d} \mathcal{N}\big(0, \sigma_{\text{EIF}}^2\big),
\]
where $\sigma_{\text{EIF}}^2 = \sigma_Y^2(1 - \rho_{Yf}^2)$ is the efficient variance.

\paragraph{Power formula.}
The power of the two-sided Wald test is
\begin{equation}\label{eq:eif-power}
\text{Power} = \Phi\left(-z_{1-\alpha/2} + \frac{|\Delta|}{\sigma_{\text{EIF}}/\sqrt{n}}\right)
+ \Phi\left(-z_{1-\alpha/2} - \frac{|\Delta|}{\sigma_{\text{EIF}}/\sqrt{n}}\right).
\end{equation}

\paragraph{Required sample size.}
Inverting~\eqref{eq:eif-power} for target power $1 - \beta$ gives
\begin{equation}\label{eq:eif-n-required}
n \ge \frac{\sigma_Y^2(1 - \rho_{Yf}^2)}{\big(\Delta / (z_{1-\alpha/2} + z_{1-\beta})\big)^2}.
\end{equation}

\paragraph{Efficiency gain.}
Compared to the classical sample mean (without predictions), which requires
\[
n_{\text{classical}} \ge \frac{\sigma_Y^2}{\big(\Delta / (z_{1-\alpha/2} + z_{1-\beta})\big)^2},
\]
the EIF-based approach reduces the required sample size by a factor of $(1 - \rho_{Yf}^2)$.
When $\rho_{Yf} = 0.9$, this is a $81\%$ reduction.

\subsection{EIF for Difference in Means}

For the two-sample setting with independent groups $A$ and $B$, the parameter of interest is
$\Delta^\star = \theta_A^\star - \theta_B^\star$.

\paragraph{Augmented influence function.}
The efficient influence function for the difference is
\[
\psi_{\lambda_A, \lambda_B}(Y_A, f_A, Y_B, f_B)
= \big(Y_A + \lambda_A(f_A - \E[f_A])\big) - \big(Y_B + \lambda_B(f_B - \E[f_B])\big).
\]

With group-specific optimal tuning:
\begin{equation}\label{eq:eif-ttest-lambda}
\lambda_A^\star = -\frac{\Cov(Y_A, f_A)}{\sigma_{f,A}^2}, \qquad
\lambda_B^\star = -\frac{\Cov(Y_B, f_B)}{\sigma_{f,B}^2},
\end{equation}
the efficient variance of the difference estimator is
\begin{equation}\label{eq:eif-ttest-var}
\Var(\widehat{\Delta}_{\text{EIF}})
= \frac{\sigma_{Y,A}^2(1 - \rho_A^2)}{n_A} + \frac{\sigma_{Y,B}^2(1 - \rho_B^2)}{n_B},
\end{equation}
where $\rho_A, \rho_B$ are the group-specific correlations between $Y$ and $f$.

\paragraph{Power formula for two-sample EIF test.}
Under balanced design ($n_A = n_B = n$) and equal variances ($\sigma_{Y,A} = \sigma_{Y,B} = \sigma_Y$),
the power is
\begin{equation}\label{eq:eif-ttest-power}
\text{Power} = \Phi\left(-z_{1-\alpha/2} + \frac{|\Delta|}{\sqrt{2}\sigma_Y\sqrt{(1-\bar\rho^2)/n}}\right)
+ \Phi\left(-z_{1-\alpha/2} - \frac{|\Delta|}{\sqrt{2}\sigma_Y\sqrt{(1-\bar\rho^2)/n}}\right),
\end{equation}
where $\bar\rho^2 = (\rho_A^2 + \rho_B^2)/2$ is the average squared correlation.

\paragraph{Required sample size per group.}
\begin{equation}\label{eq:eif-ttest-n}
n \ge \frac{2\sigma_Y^2(1 - \bar\rho^2)}{\big(\Delta / (z_{1-\alpha/2} + z_{1-\beta})\big)^2}.
\end{equation}

\subsection{Practical Calibration}

The EIF-based formulas require knowledge of:
\begin{itemize}
    \item $\sigma_Y^2$: estimated from labeled data as $\widehat{\sigma}_Y^2 = \frac{1}{n-1}\sum_{i=1}^n(Y_i - \bar{Y})^2$
    \item $\rho_{Yf}$: estimated from labeled data as the sample correlation between $Y_i$ and $f_i$
\end{itemize}

\paragraph{From MSE to correlation.}
For a well-calibrated predictor with $\E[f(X)] \approx \E[Y]$, the MSE relates to the correlation via
\[
\text{MSE} = \E[(Y - f)^2] = \sigma_Y^2 + \sigma_f^2 - 2\Cov(Y, f) = \sigma_Y^2(1 - \rho^2) + (\sigma_f - \rho\sigma_Y)^2.
\]
When $f$ is the best linear predictor of $Y$, we have $\sigma_f = \rho\sigma_Y$ and $\text{MSE} = \sigma_Y^2(1 - \rho^2)$.

\paragraph{From $R^2$ to sample size reduction.}
In regression settings, the $R^2$ of predicting $Y$ from $f$ equals $\rho_{Yf}^2$.
Thus, the sample size reduction factor is exactly $(1 - R^2)$.

\begin{tcolorbox}
\textbf{Rule of thumb:} An auxiliary predictor with $R^2 = 0.5$ (explaining 50\% of variance)
reduces the required labeled sample size by 50\%.
\end{tcolorbox}
 in main.tex when ready

\section{EIF-Based Power Analysis}
\label{sec:eif-power}

We derive power formulas from the perspective of semiparametric efficiency theory.
The efficient influence function (EIF) characterizes the asymptotically optimal
estimator for a given statistical model, and its variance directly determines
the power of hypothesis tests.

\subsection{EIF for the Mean with Auxiliary Predictions}

Consider estimating the population mean $\theta^\star = \E[Y]$ using:
\begin{itemize}
    \item Labeled data: $(X_i, Y_i, f_i)$ for $i = 1, \ldots, n$, where $f_i = f(X_i)$
    \item Unlabeled data: $(\widetilde{X}_j, \widetilde{f}_j)$ for $j = 1, \ldots, N$
\end{itemize}

\paragraph{The classical influence function.}
Without auxiliary predictions, the influence function for the mean is simply
\[
\psi_0(Y) = Y - \theta^\star.
\]
The sample mean $\bar{Y}_n$ is the corresponding M-estimator with asymptotic variance $\Var(Y)/n$.

\paragraph{The augmented influence function.}
With auxiliary predictions $f(X)$, we can construct an augmented influence function:
\begin{equation}\label{eq:eif-augmented}
\psi_\lambda(Y, f) = Y - \theta^\star + \lambda\big(f - \E[f]\big),
\end{equation}
where $\lambda \in \mathbb{R}$ is a tuning parameter.

The key observation is that $\E[\psi_\lambda(Y, f)] = 0$ for any $\lambda$ since
$\E[Y] = \theta^\star$ and $\E[f - \E[f]] = 0$. This means the estimating equation
\[
\frac{1}{n}\sum_{i=1}^n \psi_\lambda(Y_i, f_i) = 0
\]
yields an unbiased estimator of $\theta^\star$ for any choice of $\lambda$.

\paragraph{Variance of the augmented estimator.}
The variance of $\psi_\lambda(Y, f)$ is
\begin{align*}
\Var\big(\psi_\lambda(Y, f)\big)
&= \Var(Y) + \lambda^2 \Var(f) + 2\lambda\Cov(Y - \theta^\star, f - \E[f])\\
&= \sigma_Y^2 + \lambda^2 \sigma_f^2 + 2\lambda\Cov(Y, f).
\end{align*}

Minimizing over $\lambda$ gives the \emph{efficient} choice:
\begin{equation}\label{eq:eif-lambda-star}
\lambda^\star_{\text{EIF}} = -\frac{\Cov(Y, f)}{\sigma_f^2},
\end{equation}
with optimal variance
\begin{equation}\label{eq:eif-opt-var}
\Var\big(\psi_{\lambda^\star}(Y, f)\big)
= \sigma_Y^2 - \frac{\Cov(Y, f)^2}{\sigma_f^2}
= \sigma_Y^2 \big(1 - \rho^2_{Yf}\big),
\end{equation}
where $\rho_{Yf} = \Cov(Y, f)/(\sigma_Y \sigma_f)$ is the correlation between $Y$ and $f$.

\paragraph{Connection to PPI++.}
The PPI++ estimator~\eqref{eq:ppi++-mean-est} with optimal $\lambda^\star$ from~\eqref{eq:lambda-star-mean}
achieves exactly this efficient variance in~\eqref{eq:eif-opt-var}, confirming that
PPI++ is the semiparametrically efficient estimator for the mean when auxiliary predictions are available.

\subsection{Power Inversion for EIF-Based Inference}

Given the efficient variance~\eqref{eq:eif-opt-var}, we can derive power formulas by
the standard approach of inverting the power equation.

\paragraph{Setup.}
Consider testing
\[
H_0: \theta^\star = \theta_0 \quad \text{vs} \quad H_1: \theta^\star = \theta_0 + \Delta
\]
at significance level $\alpha$ with target power $1 - \beta$.

\paragraph{Asymptotic distribution under $H_1$.}
The efficient estimator $\widehat{\theta}_{\text{EIF}}$ satisfies
\[
\sqrt{n}\big(\widehat{\theta}_{\text{EIF}} - \theta^\star\big) \xrightarrow{d} \mathcal{N}\big(0, \sigma_{\text{EIF}}^2\big),
\]
where $\sigma_{\text{EIF}}^2 = \sigma_Y^2(1 - \rho_{Yf}^2)$ is the efficient variance.

\paragraph{Power formula.}
The power of the two-sided Wald test is
\begin{equation}\label{eq:eif-power}
\text{Power} = \Phi\left(-z_{1-\alpha/2} + \frac{|\Delta|}{\sigma_{\text{EIF}}/\sqrt{n}}\right)
+ \Phi\left(-z_{1-\alpha/2} - \frac{|\Delta|}{\sigma_{\text{EIF}}/\sqrt{n}}\right).
\end{equation}

\paragraph{Required sample size.}
Inverting~\eqref{eq:eif-power} for target power $1 - \beta$ gives
\begin{equation}\label{eq:eif-n-required}
n \ge \frac{\sigma_Y^2(1 - \rho_{Yf}^2)}{\big(\Delta / (z_{1-\alpha/2} + z_{1-\beta})\big)^2}.
\end{equation}

\paragraph{Efficiency gain.}
Compared to the classical sample mean (without predictions), which requires
\[
n_{\text{classical}} \ge \frac{\sigma_Y^2}{\big(\Delta / (z_{1-\alpha/2} + z_{1-\beta})\big)^2},
\]
the EIF-based approach reduces the required sample size by a factor of $(1 - \rho_{Yf}^2)$.
When $\rho_{Yf} = 0.9$, this is a $81\%$ reduction.

\subsection{EIF for Difference in Means}

For the two-sample setting with independent groups $A$ and $B$, the parameter of interest is
$\Delta^\star = \theta_A^\star - \theta_B^\star$.

\paragraph{Augmented influence function.}
The efficient influence function for the difference is
\[
\psi_{\lambda_A, \lambda_B}(Y_A, f_A, Y_B, f_B)
= \big(Y_A + \lambda_A(f_A - \E[f_A])\big) - \big(Y_B + \lambda_B(f_B - \E[f_B])\big).
\]

With group-specific optimal tuning:
\begin{equation}\label{eq:eif-ttest-lambda}
\lambda_A^\star = -\frac{\Cov(Y_A, f_A)}{\sigma_{f,A}^2}, \qquad
\lambda_B^\star = -\frac{\Cov(Y_B, f_B)}{\sigma_{f,B}^2},
\end{equation}
the efficient variance of the difference estimator is
\begin{equation}\label{eq:eif-ttest-var}
\Var(\widehat{\Delta}_{\text{EIF}})
= \frac{\sigma_{Y,A}^2(1 - \rho_A^2)}{n_A} + \frac{\sigma_{Y,B}^2(1 - \rho_B^2)}{n_B},
\end{equation}
where $\rho_A, \rho_B$ are the group-specific correlations between $Y$ and $f$.

\paragraph{Power formula for two-sample EIF test.}
Under balanced design ($n_A = n_B = n$) and equal variances ($\sigma_{Y,A} = \sigma_{Y,B} = \sigma_Y$),
the power is
\begin{equation}\label{eq:eif-ttest-power}
\text{Power} = \Phi\left(-z_{1-\alpha/2} + \frac{|\Delta|}{\sqrt{2}\sigma_Y\sqrt{(1-\bar\rho^2)/n}}\right)
+ \Phi\left(-z_{1-\alpha/2} - \frac{|\Delta|}{\sqrt{2}\sigma_Y\sqrt{(1-\bar\rho^2)/n}}\right),
\end{equation}
where $\bar\rho^2 = (\rho_A^2 + \rho_B^2)/2$ is the average squared correlation.

\paragraph{Required sample size per group.}
\begin{equation}\label{eq:eif-ttest-n}
n \ge \frac{2\sigma_Y^2(1 - \bar\rho^2)}{\big(\Delta / (z_{1-\alpha/2} + z_{1-\beta})\big)^2}.
\end{equation}

\subsection{Practical Calibration}

The EIF-based formulas require knowledge of:
\begin{itemize}
    \item $\sigma_Y^2$: estimated from labeled data as $\widehat{\sigma}_Y^2 = \frac{1}{n-1}\sum_{i=1}^n(Y_i - \bar{Y})^2$
    \item $\rho_{Yf}$: estimated from labeled data as the sample correlation between $Y_i$ and $f_i$
\end{itemize}

\paragraph{From MSE to correlation.}
For a well-calibrated predictor with $\E[f(X)] \approx \E[Y]$, the MSE relates to the correlation via
\[
\text{MSE} = \E[(Y - f)^2] = \sigma_Y^2 + \sigma_f^2 - 2\Cov(Y, f) = \sigma_Y^2(1 - \rho^2) + (\sigma_f - \rho\sigma_Y)^2.
\]
When $f$ is the best linear predictor of $Y$, we have $\sigma_f = \rho\sigma_Y$ and $\text{MSE} = \sigma_Y^2(1 - \rho^2)$.

\paragraph{From $R^2$ to sample size reduction.}
In regression settings, the $R^2$ of predicting $Y$ from $f$ equals $\rho_{Yf}^2$.
Thus, the sample size reduction factor is exactly $(1 - R^2)$.

\begin{tcolorbox}
\textbf{Rule of thumb:} An auxiliary predictor with $R^2 = 0.5$ (explaining 50\% of variance)
reduces the required labeled sample size by 50\%.
\end{tcolorbox}
